\documentclass[11pt]{article}

\usepackage[margin=1.1in]{geometry}
\usepackage{amsmath,amssymb,amsthm}
\usepackage{booktabs}
\usepackage{array}
\usepackage[colorlinks=true,linkcolor=blue,citecolor=blue,urlcolor=blue]{hyperref}

\newtheorem{theorem}{Theorem}[section]
\newtheorem{lemma}[theorem]{Lemma}
\newtheorem{corollary}[theorem]{Corollary}
\newtheorem{remark}[theorem]{Remark}
\newtheorem{definition}[theorem]{Definition}

\newcommand{\ZZ}{\mathbb{Z}}
\newcommand{\GBox}{\mathcal{B}}

\title{Fixed-length self-avoiding walks in grid boxes are eventually
polynomial:\\ proofs of twenty conjectured formulas from the OEIS}

\author{Chen-Ting Lin\thanks{\texttt{chenting@uchicago.edu}}}
\date{August 2026 (draft)}

\begin{document}
\maketitle

\begin{abstract}
For fixed $k$ and $d$, let $a_{k,d}(n)$ denote the number of $k$-step
self-avoiding walks in the $d$-dimensional grid box
$\{0,\dots,n-1\}^d$, summed over all starting cells, where---following
R.~H.~Hardin's convention in the On-Line Encyclopedia of Integer Sequences
(OEIS)---a $k$-step walk is a sequence of $k$ distinct pairwise-adjacent
cells and each direction is counted separately.  We prove that
$a_{k,d}(n)$ agrees with a polynomial in $n$ of degree $d$ for all
$n \ge 2k-1$, with coefficients determined by finitely many walk
enumerations in boxes of side at most $2k-1$.  Combining this with exact
computations that close the finitely many cases below the threshold, we
prove the conjectured (``empirical'') formulas recorded in twenty OEIS
entries: A188148--A188155 ($d=2$, $3 \le k \le 10$), A187164--A187170
($d=3$, $3 \le k \le 9$), and A188785--A188789 ($d=4$, $2 \le k \le 6$),
conjectured by R.~H.~Hardin in 2011 and open since.  All computations are
reproducible, and key numerical certificates are formally verified in the
Lean~4 proof assistant.
\end{abstract}

\section{Introduction}

The On-Line Encyclopedia of Integer Sequences (OEIS) \cite{OEIS} contains
several thousand entries whose formulas are marked \emph{empirical} or
\emph{conjectured}: closed forms, recurrences, or generating functions
found by fitting the published terms, awaiting proof.  A substantial
family of such entries is due to R.~H.~Hardin, who in the early 2010s
contributed extensive tables counting combinatorial configurations in
finite grids, together with empirically observed formulas.  Proving
Hardin's conjectures has recently become an active pastime:
Dougherty-Bliss and Kauers \cite{DBK} proved his conjectures on
``Hardinian arrays'' via the transfer-matrix method, the key structural
result being that the counts are eventually polynomial in the grid
dimension; and individual entries have been settled directly on the OEIS,
e.g.\ Corneth's 2023 proof note in entry A186864 (5-step king's tours).

This paper settles one such family in a uniform way.  For fixed $k \ge 1$
and $d \ge 1$ consider the box
$\GBox_n^d = \{0,1,\dots,n-1\}^d \subseteq \ZZ^d$, with cells adjacent when
they differ by $\pm 1$ in exactly one coordinate, and let
\[
  a_{k,d}(n) \;=\; \#\bigl\{\, k\text{-step self-avoiding walks in }
  \GBox_n^d \,\bigr\},
\]
summed over all starting cells.  Here a \emph{$k$-step walk} is a sequence
of $k$ \emph{distinct} pairwise-adjacent cells $c_1,\dots,c_k$ (thus
traversing $k-1$ edges), and a walk and its reversal are counted
separately.  This is Hardin's convention, which we confirmed against the
example section of OEIS entry A188148; we adopt it so that our statements
match the OEIS entries verbatim.

Hardin computed $a_{k,d}(n)$ for the twenty pairs $(k,d)$ listed in
Table~\ref{tab:results} and conjectured, in each case, that $a_{k,d}(n)$
is given by an explicit polynomial of degree $d$ for all $n$ beyond a
small explicit threshold $n_0(k,d)$.  These conjectures (and the
equivalent conjectured generating functions and linear recurrences added
by C.~Barker in 2018) have remained open.  We prove all of them.

\begin{theorem}[Main theorem, informal]\label{thm:informal}
For every fixed $k$ and $d$, the count $a_{k,d}(n)$ agrees for all
$n \ge 2k-1$ with a polynomial in $n$ of degree $d$ whose coefficients
are determined by finitely many walk enumerations in boxes of side at
most $2k-1$.  For the twenty pairs $(k,d)$ of
Table~\ref{tab:results}, the resulting polynomials coincide with the
conjectured OEIS formulas, and direct computation for the finitely many
$n$ between the conjectured threshold and $2k-1$ completes the proof of
each conjecture exactly as stated.
\end{theorem}

The structural part (Section~\ref{sec:theorem}) is an elementary locality
argument: a walk of $m = k-1$ edges started at a cell $c$ only sees the
box within $L^\infty$-distance $m$ of $c$, so the number of walks from
$c$ depends only on the \emph{clipped profile} of $c$---its distance to
each face of the box, capped at $m$.  Grouping cells by profile
immediately yields eventual polynomiality.  We do not claim this
observation is deep; its value lies in the uniform statement with an
explicit threshold, which reduces each OEIS conjecture to a finite,
machine-checkable computation, and in the resulting complete and
independently verifiable proofs of twenty long-standing conjectured
formulas.  This mirrors the situation in \cite{DBK}, where eventual
polynomiality (there, via transfer matrices) was likewise the engine
turning empirical formulas into theorems.

\paragraph{Verification and formalization.}
All computations were performed by two independent methods --- direct
enumeration and the profile method --- which agree on an extensive common
test set, and the profile method reproduces every term published in the
OEIS entries (22--43 terms per entry).  The polynomial coefficients are recovered by exact rational
interpolation.  In addition, all twenty conjectures have been formalized and
machine-checked end to end in the Lean~4 proof assistant
(Remark~\ref{rem:lean}).  Code and Lean sources are available at
\url{https://github.com/d26521082/saw-polynomials}.

\section{Definitions and conventions}\label{sec:defs}

Fix $d \ge 1$.  Cells are points of $\ZZ^d$; two cells are
\emph{adjacent} when they differ by $\pm 1$ in exactly one coordinate.
For $n \ge 1$ let $\GBox_n^d = \{0,\dots,n-1\}^d$.

\begin{definition}[Walks; Hardin's convention]
A \emph{$k$-step self-avoiding walk} in $\GBox_n^d$ is a tuple
$(c_1,\dots,c_k)$ of distinct cells of $\GBox_n^d$ with $c_i$ adjacent to
$c_{i+1}$ for $1 \le i < k$.  It traverses $m := k-1$ edges.  Reversals
are counted as distinct walks.  We write $a_{k,d}(n)$ for the number of
such walks, and, for a cell $c$, $f_n(c)$ for the number of walks with
$c_1 = c$, so that $a_{k,d}(n) = \sum_{c \in \GBox_n^d} f_n(c)$.
\end{definition}

The term ``$k$-step'' for a walk on $k$ cells follows the OEIS entries;
the example section of A188148 (``3-step'' walks occupying three cells)
fixes the convention unambiguously.  Note $a_{1,d}(n) = n^d$ and
$a_{2,d}(n) = 2d\,n^{d-1}(n-1)$ (each direction of each edge).

\section{The eventual polynomiality theorem}\label{sec:theorem}

Throughout this section $k, d$ are fixed and $m = k - 1$.

\begin{lemma}[Locality]\label{lem:locality}
For $c \in \GBox_n^d$ define the \emph{clipped profile}
\[
  \pi(c) \;=\; \bigl(\min(c^{(i)}, m),\ \min(n-1-c^{(i)}, m)\bigr)_{i=1}^{d}
  \;\in\; \bigl(\{0,\dots,m\}^2\bigr)^d .
\]
Then $f_n(c)$ depends only on $\pi(c)$; that is, there is a function $w$
on profiles, independent of $n$, with $f_n(c) = w(\pi(c))$ for all $n$
and all $c \in \GBox_n^d$.
\end{lemma}

\begin{proof}
A walk of $m$ edges from $c$ changes each coordinate by at most $1$ per
step, hence stays inside
$R(c) = \prod_{i=1}^{d} [\,c^{(i)}-\ell_i,\ c^{(i)}+r_i\,]$
where $\ell_i = \min(c^{(i)}, m)$ and $r_i = \min(n-1-c^{(i)}, m)$:
indeed a cell of $\GBox_n^d$ within $m$ steps of $c$ along axis $i$ lies
between $\max(0, c^{(i)}-m) = c^{(i)}-\ell_i$ and
$\min(n-1, c^{(i)}+m) = c^{(i)}+r_i$.  Conversely every cell of $R(c)$
lies in $\GBox_n^d$.  Thus the walks from $c$ in $\GBox_n^d$ are exactly
the walks from $c$ in the box $R(c)$, and translating $R(c)$ by $-c +
(\ell_1,\dots,\ell_d)$ gives a bijection with walks in
$\prod_i \{0,\dots,\ell_i+r_i\}$ started at $(\ell_1,\dots,\ell_d)$.
The latter count depends only on
$(\ell_i, r_i)_i = \pi(c)$.
\end{proof}

Each value $w(p)$ is a finite, explicitly computable integer: enumerate
the walks of $m$ edges in a box of side lengths $\ell_i + r_i + 1 \le
2m+1$.  By reflection and permutation of axes, $w$ is invariant under
swapping $\ell_i \leftrightarrow r_i$ and permuting the $d$ pairs, which
we exploit in computations.

\begin{lemma}[Census]\label{lem:census}
Let $n \ge 2m+1$.  Along one axis, as $x$ ranges over $\{0,\dots,n-1\}$,
the clipped pair $\bigl(\min(x,m), \min(n-1-x,m)\bigr)$ takes the value
$(i, m)$ exactly once for each $0 \le i \le m-1$ (namely $x = i$), the
value $(m, j)$ exactly once for each $0 \le j \le m-1$ (namely
$x = n-1-j$), and the value $(m, m)$ exactly $n - 2m$ times; no other
values occur.
\end{lemma}

\begin{proof}
If $x < m$ then, since $n - 1 - x > n - 1 - m \ge m$, the pair is
$(x, m)$.  Symmetrically if $x > n-1-m$ the pair is $(m, n-1-x)$ with
$n-1-x < m$.  The remaining $x \in [m, n-1-m]$, of which there are
$n-2m$, give $(m,m)$.  The three ranges are disjoint because
$m \le n-1-m$.
\end{proof}

\begin{theorem}[Eventual polynomiality]\label{thm:main}
For all $n \ge 2m + 1 = 2k - 1$,
\[
  a_{k,d}(n) \;=\; \sum_{S \subseteq \{1,\dots,d\}} (n-2m)^{|S|}
  \sum_{\substack{p = (p_1,\dots,p_d)\\ p_i = (m,m) \text{ for } i \in S\\
      p_i \in E \text{ for } i \notin S}} w(p),
\]
where $E = \{(i,m) : 0 \le i < m\} \cup \{(m,j) : 0 \le j < m\}$ is the
set of boundary pairs.  In particular $a_{k,d}(n)$ agrees on
$n \ge 2k-1$ with a polynomial in $n$ of degree (at most) $d$, whose
leading coefficient is the count $w\bigl((m,m)^d\bigr)$ of $k$-step walks
from a central cell of the $(2m{+}1)^d$ box.
\end{theorem}

\begin{proof}
By Lemma~\ref{lem:locality}, $a_{k,d}(n) = \sum_{p} N_n(p)\, w(p)$ where
$N_n(p)$ is the number of cells with profile $p$.  Profiles factor across
axes, so by Lemma~\ref{lem:census}
$N_n(p) = \prod_{i=1}^d \nu_n(p_i)$ with $\nu_n((m,m)) = n-2m$,
$\nu_n(e) = 1$ for $e \in E$, and $\nu_n = 0$ otherwise.  Collecting
terms by the set $S$ of axes carrying the interior pair $(m,m)$ gives the
displayed formula, a polynomial in $n-2m$ (hence in $n$) of degree at
most $d$.  The degree-$d$ coefficient is the $S = \{1,\dots,d\}$ term,
$w((m,m)^d) \ge 1$ (at least the straight walk exists).
\end{proof}

\begin{corollary}[Generating functions and recurrences]\label{cor:gf}
Fix $(k,d)$ and suppose $a_{k,d}(n) = P(n)$ for all $n > n_0$, with
$P$ a polynomial of degree $d$.  Then $\sum_{n\ge1} a_{k,d}(n) x^n = Q(x) + R(x)/(1-x)^{d+1}$ for
polynomials $Q, R$ with $\deg R \le d$, where $Q$ absorbs the terms
$n \le n_0$, and
$a_{k,d}(n) = \sum_{i=1}^{d+1} (-1)^{i+1}\binom{d+1}{i} a_{k,d}(n-i)$
for $n > n_0 + d + 1$.  Hence the conjectured Barker generating functions
and recurrences in the OEIS entries follow from the polynomial formulas.
\end{corollary}

\begin{proof}
Standard: the $(d{+}1)$-st finite difference of a degree-$d$ polynomial
vanishes.
\end{proof}

\section{Proofs of the twenty conjectures}\label{sec:results}

For each pair $(k,d)$ in Table~\ref{tab:results} we proceed in three
steps, all exact integer computations.

\begin{enumerate}
\item \textbf{Coefficients.}  Compute $a_{k,d}(n)$ by the profile method
  (Lemmas~\ref{lem:locality}--\ref{lem:census}: enumerate $w(p)$ for the
  finitely many profiles, then form the census sum) at $d+1$ points
  $n \ge 2k-1$, and recover the polynomial $P_{k,d}$ of
  Theorem~\ref{thm:main} by exact rational interpolation.  In all twenty
  cases $P_{k,d}$ \emph{equals the conjectured OEIS polynomial,
  coefficient by coefficient}.
\item \textbf{Gap.}  The OEIS conjectures assert validity for
  $n > n_0$ with $n_0 < 2k-1$.  For the finitely many
  $n_0 < n < 2k-1$ we evaluate $a_{k,d}(n)$ exactly (the census sum is
  valid for every $n$, using the true multiset of clipped pairs for that
  $n$) and check $a_{k,d}(n) = P_{k,d}(n)$.  All checks pass.
\item \textbf{Consistency.}  Independently, direct enumeration over all
  starting cells agrees with the profile method wherever both were run,
  and the profile method reproduces \emph{every} term published in the
  OEIS entry (22--43 terms per entry, fetched verbatim from the OEIS).
\end{enumerate}

Steps 1--2, together with Theorem~\ref{thm:main}, constitute a complete
proof of each conjectured formula; step 3 is an additional guard against
implementation error.  By Corollary~\ref{cor:gf} the conjectured
generating functions and recurrences in the same entries follow.

\begin{table}[ht]
\centering
\small
\begin{tabular}{@{}llll@{}}
\toprule
OEIS & $(k,d)$ & $a_{k,d}(n)$ for $n > n_0$ & $n_0$ \\
\midrule
A188148 & $(3,2)$  & $12n^2-24n+8$ & $1$ \\
A188149 & $(4,2)$  & $36n^2-100n+56$ & $2$ \\
A188150 & $(5,2)$  & $100n^2-360n+272$ & $3$ \\
A188151 & $(6,2)$  & $284n^2-1228n+1152$ & $4$ \\
A188152 & $(7,2)$  & $780n^2-3960n+4432$ & $5$ \\
A188153 & $(8,2)$  & $2172n^2-12500n+16096$ & $6$ \\
A188154 & $(9,2)$  & $5916n^2-38192n+55600$ & $7$ \\
A188155 & $(10,2)$ & $16268n^2-115548n+186528$ & $8$ \\
\midrule
A187164 & $(3,3)$  & $30n^3-60n^2+24n$ & $1$ \\
A187165 & $(4,3)$  & $150n^3-426n^2+312n-48$ & $2$ \\
A187166 & $(5,3)$  & $726n^3-2640n^2+2688n-720$ & $3$ \\
A187167 & $(6,3)$  & $3534n^3-15366n^2+19536n-7056$ & $4$ \\
A187168 & $(7,3)$  & $16926n^3-85380n^2+128832n-57312$ & $5$ \\
A187169 & $(8,3)$  & $81390n^3-463074n^2+801216n-418032$ & $6$ \\
A187170 & $(9,3)$  & $387966n^3-2452704n^2+4766544n-2833872$ & $7$ \\
\midrule
A188785 & $(2,4)$  & $8n^4-8n^3$ & $0$ \\
A188786 & $(3,4)$  & $56n^4-112n^3+48n^2$ & $1$ \\
A188787 & $(4,4)$  & $392n^4-1128n^3+912n^2-192n$ & $2$ \\
A188788 & $(5,4)$  & $2696n^4-9968n^3+11424n^2-4416n+384$ & $3$ \\
A188789 & $(6,4)$  & $18584n^4-82552n^3+119616n^2-64320n+9984$ & $4$ \\
\bottomrule
\end{tabular}
\caption{The twenty proved formulas.  In each case the polynomial is the
$P_{k,d}$ of Theorem~\ref{thm:main} and the threshold $n_0$ is exactly as
conjectured in the OEIS entry.}
\label{tab:results}
\end{table}

\begin{remark}[Sharpness of the thresholds]
The conjectured thresholds are exactly minimal: for every entry with
$n_0 \ge 1$ we verified by exact computation that $a_{k,d}(n_0) \ne
P_{k,d}(n_0)$ (e.g.\ $a_{10,2}(3) = 0$ while $P_{10,2}(3) = -13704$).
For small $n$ the clipped neighborhoods of opposite faces overlap and the
census of Lemma~\ref{lem:census} no longer applies.  That the true
threshold $n_0$ is smaller than our guaranteed $2k-1$ reflects
cancellation among boundary corrections, which the gap computation of
step~2 captures; we did not seek a closed-form explanation.
\end{remark}

\begin{remark}[Formal verification]\label{rem:lean}
All twenty statements have been formalized and machine-checked end to
end in Lean~4 with mathlib.  For each $d \in \{2,3,4\}$, the locality
lemma, translation invariance, the single-axis census, and the
polynomial form of Theorem~\ref{thm:main} are proved as theorems
depending only on Lean's standard axioms (\texttt{propext},
\texttt{Classical.choice}, \texttt{Quot.sound}); the twenty OEIS entry
theorems combine these with window constants and sub-threshold values
evaluated by the compiler via \texttt{native\_decide}, adding only the
standard \texttt{Lean.ofReduceBool} trust assumption.  For example, the
Lean statement of A188148 reads: for all $n \ge 2$,
$(A\,2\,n : \mathbb{Z}) = 12n^2 - 24n + 8$, where $A\,m\,n$ is the
walk count defined from first principles.  In addition, an independent
list-based Lean implementation of the count is checked by
\texttt{native\_decide} against the published OEIS data.
\end{remark}

\section{Extensions}\label{sec:extensions}

The argument uses nothing about the nearest-neighbor step set beyond
finiteness, and nothing about the box being a cube.

\begin{enumerate}
\item \emph{Other step sets.}  For any finite step set
  $S \subseteq \ZZ^d \setminus \{0\}$ (king moves, knight moves, etc.)
  with $M = \max_{s \in S} \|s\|_\infty$, the same proof shows that the
  number of $k$-step $S$-walks (self-avoiding or not, or with any local
  constraint that is a function of the walk shape) in $\GBox_n^d$ is
  polynomial of degree $d$ in $n$ for $n \ge 2mM + 1$.  This covers, for
  instance, Hardin's king-tour family A186861--A186870, one member of
  which (A186864) was proved by Corneth by a related argument.
\item \emph{Rectangular boxes.}  Replacing $\GBox_n^d$ by
  $\prod_i \{0,\dots,n_i-1\}$ yields, verbatim, a multivariate polynomial
  in $(n_1,\dots,n_d)$ once every $n_i \ge 2m+1$, proving the
  row-and-column formulas in the associated OEIS arrays (A188147,
  A187162, A188784).
\item \emph{Walks with marked endpoints, closed walks.}  Any statistic
  determined by the translation class of the walk and its clipped profile
  admits the same treatment; e.g.\ counts of self-avoiding \emph{polygons}
  by perimeter in $\GBox_n^d$ are likewise eventually polynomial.
\end{enumerate}

We intend the Lean formalization to cover the general statement in
item~1.

\section*{Acknowledgments}

The computations, drafting, and Lean development were carried out with
substantial assistance from Claude (Anthropic).  The author is
responsible for all claims.

\begin{thebibliography}{9}

\bibitem{OEIS}
OEIS Foundation Inc.\ (2026), \emph{The On-Line Encyclopedia of Integer
Sequences}, \url{https://oeis.org}.  Entries A188147--A188155,
A187162--A187170, A188784--A188789, A186861--A186870.

\bibitem{DBK}
R.~Dougherty-Bliss and M.~Kauers, \emph{Hardinian arrays},
Electron.\ J.\ Combin.\ \textbf{31}(2) (2024), \#P2.9; arXiv:2309.00487.

\bibitem{Kauers}
M.~Kauers, \emph{Some D-finite and some possibly D-finite sequences in
the OEIS}, 2023.

\bibitem{MNP}
L.~Major, L.~N\'emeth, A.~Pahikkala, and L.~Szalay,
\emph{Self-avoiding walks of specified lengths on rectangular grid
graphs}, Aequationes Math.\ \textbf{98} (2024), 215--239.

\bibitem{JCTA}
H.~Belbachir, L.~Major, L.~N\'emeth, and L.~Szalay,
\emph{Step-constrained self-avoiding walks on finite grids},
J.\ Combin.\ Theory Ser.\ A \textbf{218} (2026), 106104.

\bibitem{Lean}
L.~de~Moura and S.~Ullrich, \emph{The Lean 4 theorem prover and
programming language}, CADE 2021.

\end{thebibliography}

\end{document}
